\section{The source bridges are the integration, not a footnote}
\label{sec:bridges}

\subsection{A checked finite model is not yet a proof of the original goal}
Forge's existing design already separates source identity, search guidance,
certificate acceptance, and theorem authority \cite{forge-leant}. Those
mechanisms should be reused, not redesigned here. The new obligation is more
specific: each worker has a different \emph{semantic direction} and therefore
needs a different translation theorem.

The initial implementation should take an explicit model and a proved bridge.
Only after this route works should a reifier infer the model from arbitrary
Lean definitions. That order gives an end-to-end theorem sooner and prevents
an attractive front-end demo from concealing an unproved encoding.
A useful worker result is thus conceptually a dependent pair
\[
   \bigl(\text{model certificate},\ \text{proof of the model proposition}\bigr),
\]
which is then consumed by the source bridge. The model proposition is not
allowed to change when proof search is inconvenient.

\subsection{Reachability: forward simulation and witness lifting differ}
Let $\mathcal C$ be a source transition system, let $\mathcal A$ be the
pushdown model, and let $\mathcal R(c,a)$ relate source and abstract
configurations. To transfer an abstract safety result, it is sufficient to
prove an initial relation, preservation of the relation along every source
step by a finite abstract run, and a target mapping
\[
 \operatorname{Bad}_{\mathcal C}(c)\wedge\mathcal R(c,a)
 \Longrightarrow\operatorname{Target}_{\mathcal A}(a).
\]
An abstract nonreachability theorem then rules out a source error. Abstract
behavior may be an over-approximation; extra abstract paths weaken the ability
to prove safety but do not invalidate a proved safety result.

A positive abstract run is different. It requires a lifting theorem that
constructs a related source run, or an exact operational correspondence. A
spurious abstract path cannot establish a source witness. A practical first
adapter should use an equivalence of the two reachability propositions, with
both directions proved once for its recognized syntax.

For regular stack targets, the target DFA is also part of the bridge. The
compiler theorem must connect \emph{the original stack predicate} to the DFA
language, then to the monitor phase, then to the empty-stack query. A proof that
the compiled automaton empties its stack does not by itself identify what that
stack meant in the source program. Top-first order and the fresh bottom symbol
must occur in the theorem, not only in a comment in the exporter.

\subsection{Games need alternating, ownership-aware refinement}
A forward simulation of unlabelled graphs is not sufficient for strategy
synthesis. The source may distinguish actions the program controls from
outcomes controlled by an environment. At a source state $x$, an abstract
protagonist move must have a legal \emph{uniformly realizable} implementation;
every possible environment outcome of that implementation must be accounted
for.

A precise first interface uses a turn-based source arena with an abstraction
$\alpha:X\to V$. For a positive policy $\sigma$ on $W$, require a source
invariant $I$, the initial condition $I(x_0)$ and $\alpha(x_0)\in W$, and the
following obligations:
\begin{enumerate}
\item Ownership is preserved. At a protagonist state satisfying $I$, a
  checked realizer chooses a legal source successor mapping to
  $\sigma(\alpha(x))$ and preserving $I$.
\item At an antagonist state satisfying $I$, \emph{every} legal source
  successor preserves $I$ and maps to a successor of $\alpha(x)$ in the
  original abstract arena.
\item Abstract acceptance implies source acceptance:
  $I(x)\wedge\alpha(x)\in B\Rightarrow B_{\rm src}(x)$.
  Source plays are infinite under the relevant strategy, or the source's
  treatment of termination is explicitly reflected in the arena.
\end{enumerate}
Under these conditions the abstract policy lifts to a source policy, every
consistent source play maps stepwise to a consistent abstract play, and
Theorem~\ref{thm:buchi-rank} transfers. The realizer is a typed function whose
correctness is quantified over all source states satisfying $I$, not a lookup
that was tested on one representative of each abstract block.

When a source action has several possible outcomes, selecting the action is
not the same as selecting an outcome. A concrete encoding inserts an
intermediate antagonist-owned action vertex: the protagonist chooses the
vertex for the action, then the antagonist chooses its outcome. Labels on
these intermediate vertices and the two-step-to-one-step correspondence must
be fixed in the bridge. Simply turning every possible action outcome into a
protagonist edge would overstate the program's power.

A negative game certificate reverses the controlled side. To conclude that
\emph{no source strategy} succeeds, the abstract antagonist policy must lift,
all protagonist choices must be represented, and source accepting visits must
imply the relevant abstract accepting visits. The positive label implication
alone is not enough for this direction. The implementation should expose
separate positive and negative bridge types rather than a single vaguely named
``simulation'' field.

\subsection{Stuttering can destroy liveness even when it preserves safety}
Suppose the abstract system takes $a\to b$, with $b$ accepting forever. Suppose
the source can perform infinitely many internal steps that all abstract to
$a$. A weak simulation that ignores internal steps may still relate every
finite source prefix to the abstract initial prefix. It does not establish
that the source ever reaches an accepting state.

There are two suitable first implementations. The simpler one requires each
source step to map to an abstract step, as above. The more flexible one allows
stuttering only when supplied with a well-founded rank that strictly decreases
during internal steps and is reset only when an abstract step occurs. An
infinite source play must then contain infinitely many abstract steps. If
acceptance transfer is also aligned with those steps, abstract B\"uchi
recurrence implies source recurrence. This use of a stuttering rank complements
the new reset-at-acceptance rank; the two ranks certify different things.

Quantitative recurrence-gap bounds do not transfer through this more flexible
bridge without a bound on the number or duration of internal steps. Qualitative
liveness and bounded response time should remain distinct result types.

\subsection{Probability requires exact pushforward, not support simulation}
Let a source probabilistic program have transition distribution $K(x)$ and let
$\alpha:X\to S$ map source states to the finite chain. A sufficient exact
quotient condition is
\begin{equation}
 \alpha_*K(x)=P_{\alpha(x),\cdot}
 \qquad\text{for every source state satisfying the invariant.}
 \label{eq:lumpability}
\end{equation}
Also require invariant preservation with probability one, the initial-state
mapping, and exact correspondence of terminals, running costs, and terminal
payoffs:
\begin{align*}
 x\in T_{\rm src}&\Longleftrightarrow\alpha(x)\in T,\\
 c_{\rm src}(x)&=c(\alpha(x)),\\
 g_{\rm src}(x)&=g(\alpha(x))\quad(x\in T_{\rm src}).
\end{align*}
Induction on finite horizons then equates the abstract distributions with the
pushforwards of the source distributions. The finite-horizon identities and
limits from Section~\ref{sec:markov} transfer to the source.

This is a strong but useful condition: the source state space can be infinite
provided the finite quotient is exact. It is not enough that source and
abstract transitions have the same possible endpoints. Two source states with
exit probabilities $1/4$ and $1/2$ cannot generally be merged into one abstract
state with a single transition row. An average row has no justification merely
because it falls between the two probabilities. Formula~\eqref{eq:lumpability}
asks for mass equality for every abstract block, not an empirical average.

Mathlib's probability mass functions and probability kernels provide the
appropriate semantic objects \cite{mathlib-pmf,mathlib-kernel}. For discrete
source programs, the intended equation is a distribution-map equality; for a
measure-theoretic kernel it is the corresponding pushforward equality. The
finite rational matrix must be linked to that object with proofs of
nonnegativity and unit mass. No claim is made that importing the probability
library supplies this application-specific quotient theorem automatically.

\subsection{What to borrow from Leant and Djex, without repeating their work}
Leant and Djex already preserve source-owned typed candidates, universe and
binder selections, and scoped provider evidence, and Leant re-elaborates
candidate terms \cite{leant,djex}. The new use is to ask their synthesis layer
for the \emph{realizer} or \emph{policy implementation} demanded by an accepted
behavioral certificate.

For example, the game solver can propose a finite abstract policy while the
typed synthesizer builds a source action-selection function from available
providers. Forge then proves the universal realization and refinement
conditions. A syntactically different function of the same type is not an
acceptable substitute if its behavior differs. Conversely, a well-typed source
function does not receive the temporal theorem merely because its finite
observations agree with the policy on a sample.

The semantic proof belongs to the exact retained candidate, with its lexical
providers and dictionaries. This is an application of the neighboring
projects' existing ownership discipline, not a new discipline claimed by this
article. It also does not resolve their separate open Church-encoding or
higher-universe synthesis cases. The proposed integration adds temporal
contracts to selected typed candidates; it does not assert a complete
higher-order program synthesizer.
